\section{Chapitre~\ref{sec:dendrites}. Le nombre de dendrites}

La structure des classes et le nombre d'entrées sont mesurés par l'anatomie et par des enregistrements électriques; l'estimation~\eqref{eq:dendriteload} relève de la simple physique du condensateur, tandis que l'explication computationnelle du nombre d'entrées s'appuie sur la théorie et sur des modèles.

{\footnotesize
\setlength{\tabcolsep}{3pt}\renewcommand{\arraystretch}{1.15}
\begin{longtable}{>{\raggedright}p{0.5cm}>{\raggedright}p{4.0cm}>{\raggedright}p{5.6cm}>{\raggedright}p{2.3cm}>{\raggedright\arraybackslash}p{1.7cm}}
\toprule
N\textsuperscript{o} & Affirmation & Expérience et ce qui est mesuré & Source & Statut \\
\midrule\endfirsthead
\toprule
N\textsuperscript{o} & Affirmation & Expérience et ce qui est mesuré & Source & Statut \\
\midrule\endhead
1 & Capacité spécifique de la membrane $c_m\approx1$~µF/cm$^2$ & On prélevait la membrane du corps d'un neurone dans une pipette, on appliquait un échelon de tension et, d'après la décroissance du courant capacitif, on trouvait la capacité totale, puis on la divisait par la surface: $0.9$~µF/cm$^2$ pour trois classes de neurones & \cite{Gentet2000} & mesuré \\
2 & Temps de charge $t\approx c_mA\Delta V/I$~\eqref{eq:dendriteload}: un grand neurone a besoin de dizaines d'entrées simultanées, une seule grosse synapse suffit à un petit & Estimation par la loi de charge d'un condensateur; les nombres ($20$ et $300$~pF, $0.1$ et quelques nA) sont des ordres de grandeur & démonstration dans le chapitre & théorie \\
3 & Une seule synapse énorme délivre un courant de quelques nanoampères et amène aussitôt un petit neurone au seuil & Enregistrement simultané de la terminaison et du destinataire en tranche: courant d'une synapse de $2$--$13$~nA, chaque impulsion d'entrée provoque une réponse suprasliminaire, avec un délai d'environ $0.4$~ms à $36^\circ$C; le destinataire émet \nt{GLYC} & \cite{Borst1995,BorstSoria2012} & mesuré \\
4 & Zéro dendrite: dans les neurones sensoriels du toucher et de la douleur, l'impulsion va du capteur à la moelle épinière sans passer par le corps cellulaire & La structure (un axone qui se divise en deux près du corps) est établie par l'anatomie. Que la conduction n'exige pas l'excitabilité du corps a été montré sur un modèle validé d'un tel neurone: sans corps excitable, l'impulsion franchit la bifurcation avec la même fiabilité & \cite{AmirDevor2003} & théorie \\
5 & Une dendrite: le neurone olfactif porte un seul type de récepteur et transmet une seule ligne d'entrée & Revue d'expériences sur les gènes des récepteurs: chaque neurone olfactif exprime un seul gène de récepteur parmi une grande famille & \cite{Mombaerts2004} & mesuré \\
6 & Deux dendrites: dans les détecteurs de direction du son, les entrées des oreilles gauche et droite se posent sur des dendrites différentes & Anatomie et enregistrements chez les mammifères, réunis dans une revue: les entrées des deux oreilles sont réparties sur deux dendrites opposées & \cite{Grothe2010} & mesuré \\
7 & Répartir les entrées sur deux dendrites rend le \og ET\fg{} plus strict & Montré sur un modèle: la saturation~\eqref{eq:shunt} sur une seule dendrite empêche les entrées d'un même côté d'atteindre le seuil, et la détection des coïncidences s'améliore. Il n'existe pas d'expérience directe où l'on aurait modifié la disposition des entrées & \cite{AgmonSnir1998} & théorie \\
8 & Environ $7\cdot10^{10}$ petits neurones \nt{GLUT} dans le circuit d'apprentissage moteur, avec $\sim4$ entrées chacun & Comptage des cellules par fractionnement isotrope: environ $69$~milliards de neurones dans cette partie du cerveau humain. Enregistrement d'une telle cellule chez le rat vivant: un contact déclenche des bouffées de courants discrets venant de quelques entrées, et la cellule décharge quand ils s'additionnent & \cite{Azevedo2009,Chadderton2004} & mesuré \\
9 & Un élément à seuil à $n$ entrées sépare au plus $P_{\max}\approx2n$ motifs aléatoires~\eqref{eq:cover}; pour le neurone de sortie du circuit d'apprentissage moteur, la borne est $\sim2\cdot10^5$, et l'estimation réaliste $\sim4\cdot10^4$ associations & La borne est un résultat classique de géométrie combinatoire. L'estimation réaliste vient d'une théorie de la capacité avec des poids uniquement positifs et une marge contre le bruit, appliquée au nombre mesuré d'entrées de ce neurone & \cite{Cover1965,Brunel2004} & théorie \\
10 & Les mélangeurs à peu d'entrées servent à déployer l'entrée; \og quatre\fg{} est proche de l'optimum & Théorie: la dimension de la représentation que fournit une couche de mélangeurs est maximale à peu près pour le nombre d'entrées observé en anatomie; l'hypothèse initiale de l'expansion figure dans des travaux classiques & \cite{Marr1969,Albus1971,LitwinKumar2017} & hypothèse \\
11 & Grands arbres: $10^4$--$10^5$ entrées & Comptage des synapses sur des images de microscopie électronique: $\approx1.7\cdot10^5$ synapses sur un seul neurone \nt{GABA} de sortie du circuit d'apprentissage moteur chez le rat; environ $30\,000$ entrées excitatrices et $1\,700$ inhibitrices sur un neurone \nt{GLUT} de l'hippocampe & \cite{NapperHarvey1988,Megias2001} & mesuré \\
12 & Les branches d'un grand arbre sont des sous-circuits distincts dotés de leurs propres seuils; le neurone est un petit réseau multicouche & Stimulation ponctuelle de deux sites sur de fines dendrites: les entrées d'une même branche s'additionnent selon une sigmoïde, celles de branches différentes linéairement. Dans les dendrites de neurones du cortex humain, on a trouvé des impulsions calciques graduées qui permettent à une seule cellule de résoudre un problème non linéairement séparable. Modèle: pour reproduire un seul neurone, il faut un réseau profond & \cite{Polsky2004,Gidon2020,Beniaguev2021} & mesuré \\
13 & Dendrites-sorties: chaque branche d'un neurone \nt{GABA} de la rétine est un détecteur de direction distinct & Imagerie calcique à deux photons dans des branches isolées: le signal d'une branche est sélectif au mouvement allant du corps cellulaire vers la pointe, alors que la tension du corps ne l'est pas & \cite{Euler2002} & mesuré \\
14 & Le nombre d'entrées suit la tâche (proposition~\ref{prop:dendrites}) & La structure des classes est mesurée (lignes 3--13); que le nombre d'entrées soit choisi pour la vitesse ou pour la classification n'a été montré que par la théorie et des modèles, pour certaines classes & \cite{LitwinKumar2017,AgmonSnir1998} & hypothèse \\
\bottomrule
\end{longtable}}
